\begin{textsample}{NEG Dim 2 – vem\_ed\_12 – Score -1.00 – vem\_ed\_12/t053.txt}  \label{ex:f2_neg_025}
COLABORAÇÕES COM UNIVERSIDADE DE AVEIRO No segundo semestre de 2021, as Fatecs Lins, Ipiranga e Guaratinguetá iniciaram Projetos Colaborativos Internacionais com a Universidade de Aveiro (UA / Portugal).
As atividades se mantêm em 2022.
Esses intercâmbios virtuais, também chamados de COIL (Collaborative Online International Learning) ou PCIs (Projetos Colaborativos Internacionais, no Centro Paula Souza), são frutos de uma reunião realizada pela equipe dos PCIs / Cesu em 29 de julho de 2021 com diretores de curso (coordenadores) e com Ana Balula, da Equipa de Apoio para a Inovação Curricular e Pedagógica da instituição portuguesa.
Na Fatec Guaratinguetá, Sérgio Tenório dos Santos Neto, professor de Comunicação Empresarial, desenvolveu um projeto voltado a técnicas para criar um pitch de sucesso com a professora Dina Baptista, de Técnicas de Expressão Oral e Escrita (UA).
Mais detalhes sobre o PCI / COIL, que deve ser continuado no segundo semestre de 2022, estão disponíveis em: https://www.ua.pt/pt/noticias/0/72337 Na Fatec Lins Na Fatec Lins, Silvio Ribeiro, que leciona Metodologia de Projeto de Produto, e Florbela Machado, docente de Gestão da Qualidade no Turismo na UA, desenvolveram com seus estudantes um guia prático para intercâmbio de turismo acessível, no segundo semestre de 2021.
Metade dos grupos abordou o destino turístico Foz do Iguaçu (Brasil) e os demais apresentaram Aveiro (Portugal).
Os trabalhos, apresentados na ferramenta Prezi e em vídeos, foram publicados na plataforma Padlet.
Ao longo de 2022, Ribeiro fará dupla com o professor Fernando Florim de Lemos, para elaboração de um guia de turismo eletrônico Brasil / Portugal.
Neste primeiro semestre, os grupos vão levantar informações junto aos hotéis e no segundo semestre, elaborar o roteiro eletrônico.
Na Fatec Ipiranga Na Fatec Ipiranga, Patrícia Sales Patrício, da equipe dos PCIs / Cesu e professora de Comunicação Interna e Endomarketing, realizou, em novembro e dezembro de 2021, um projeto com Ana Torres, professora de Fundamentos de Marketing da UA.
Oito equipes mistas (ou equipas, como se diz em Portugal) estudaram como empresas brasileiras em Portugal (e o Grupo Sonae, português, com presença no Brasil) comunicam suas práticas de Responsabilidade Social Corporativa.
Os \textbf{pôsteres} foram publicados na plataforma Padlet e as apresentações dos grupos ocorreram no Zoom.
O projeto prossegue, com algumas alterações, no primeiro semestre de 2022.
Da UA, participam as docentes Ana Torres e Paula Pinheiro – ambas ministram a disciplina Comportamento de Consumo.
Além do estudo da comunicação das práticas de Responsabilidade Socioambiental, estudantes brasileiros e portugueses vão abordar as implicações desses processos no comportamento de consumo sustentável.

% matched lemmas: pôster
\end{textsample}
